\section{Related work}\label{sec:related}

\subsection{Spatial primitives and point-cloud libraries}

Spatial indexing methods decompose geometric space through fundamentally distinct partitioning principles. Data-driven binary space partitioning is exemplified by the kd-tree of Bentley~\cite{bentley1975kdtree}, extended to nearest-neighbor search by Friedman et~al.~\cite{friedman1977kdtree}. In point clouds, the choice of split heuristic alters the geometric properties of the resulting cells: while classical median splits guarantee balanced tree height, they frequently generate elongated, high-aspect-ratio bounding boxes in regions with non-uniform point density. The sliding-midpoint rule of Maneewongvatana and Mount~\cite{maneewongvatana1999sliding} mitigates this degeneration by bounding cell aspect ratios, ensuring that empty or sparse regions do not distort spatial pruning. Conversely, space-driven subdivisions partition coordinate domains into regular geometric cells: the quadtree of Finkel and Bentley~\cite{finkel1974quadtrees} and the octree of Meagher~\cite{meagher1980octree}, surveyed exhaustively by Samet~\cite{samet2006foundations}, use bitwise coordinate interleaving for fast cell location and hierarchical level-of-detail extraction. However, because their subdivision geometry is decoupled from data density, uniform octrees can exhibit deep, sparse branch paths when points lie on thin two-dimensional surfaces or terrain manifolds.

Object-driven bounding volume hierarchies group clusters of geometric entities under tight bounding volumes. Classical spatial database indexes include the R-tree of Guttman~\cite{guttman1984rtree}, the R$^*$-tree of Beckmann et~al.~\cite{beckmann1990rstar}, and the bulk-loaded sort-tile-recursive (STR) tree of Leutenegger et~al.~\cite{leutenegger1997str}. While effective for low-dimensional spatial ranges, overlapping bounding boxes at higher tree depths can degrade traversal efficiency in dense geometric point clouds. For parallel and graphics-oriented workloads, the linear bounding-volume hierarchy (LBVH) of Karras~\cite{karras2012lbvh} constructs hierarchies from Morton-coded primitives to maximize thread-level parallelism.

In practice, geometric processing workflows rarely implement spatial trees from first principles; instead, practitioners interact with them through established libraries. Notable examples include FLANN~\cite{muja2009flann} and its lightweight header-only derivative nanoflann~\cite{blanco2014nanoflann}, the Point Cloud Library (PCL)~\cite{rusu2011pcl}, Open3D~\cite{zhou2018open3d}, the radius-search octree of Behley et~al.~\cite{behley2015octree}, and Boost.Geometry's R-tree~\cite{boost_geometry}. Crucially, each library couples its chosen primitive to specific structural defaults---such as a fixed leaf capacity (typically 10 to 64 points) and an arbitrary maximum tree depth---alongside library-specific node and memory representation choices. As Section~\ref{sec:libraries} demonstrates, these implementation choices often introduce performance discrepancies that overshadow the theoretical characteristics of the underlying spatial primitives.

\subsection{Empirical evaluations and benchmark studies}

A recurring theme in the spatial indexing literature is the substantial gap between theoretical algorithmic complexity and practical execution speed. In an early empirical study on shape registration, Elseberg et~al.~\cite{elseberg2012comparison} compared multiple nearest-neighbor search strategies and implementations across terrestrial scans. They discovered that low-level implementation details---specifically compact node memory layouts and continuous coordinate buffering---accounted for more than a factor of two in execution throughput, frequently outweighing algorithmic distinctions between kd-trees and octrees. In a targeted study on radius search, Behley et~al.~\cite{behley2015octree} evaluated octree bucket capacities across multiple datasets, finding that leaf sizes between 32 and 64 points optimized the trade-off between tree descent overhead and linear leaf distance calculations; however, their evaluation fixed bucket parameters globally rather than tuning per dataset or query distribution.

Recent benchmarking efforts have broadened this empirical inquiry. Laso and Yermo's Cheesemap~\cite{laso2025cheesemap} benchmarked library implementations across ten distinct LiDAR point clouds, demonstrating that the optimal primitive depends heavily on the sensor acquisition modality: flat grid-based representations dominated airborne LiDAR scans, whereas hierarchical octrees proved superior on mobile and terrestrial scans. Vi\~nambres et~al.~\cite{vinambres2026linearoctrees} evaluated linear octrees combined with space-filling curves (Morton and Hilbert orderings), reporting query speedups of 10\% to 50\% attributable entirely to improved cache locality and sequential memory layout at a fixed leaf capacity. Indexicon~\cite{indexicon2026} surveyed library defaults across geographic datasets via global parameter sweeps. In a broader evaluation of learned indexing, Liu et~al.~\cite{liu2024learnedspatial} evaluated multi-dimensional learned index structures against conventional alternatives, concluding that learned approaches fail to surpass properly implemented kd-trees and STR trees on nearest-neighbor workloads.

While these empirical studies provide valuable insights into specific facets of point-cloud search, none systematically decomposes the individual sources of performance gain---from library implementation wrappers and primitive selection to leaf capacity, split policies, memory page sizes and storage placement---within an identical codebase and layout engine across scales from one million to one billion points. We perform this decomposition in Section~\ref{sec:levers}.

\subsection{Workload-aware, synthesized, and learned data layouts}

In the relational database systems community, physical index design has long been treated as an adaptive optimization problem governed by query workloads. Foundational work in automated index selection (e.g., Chaudhuri and Narasayya~\cite{chaudhuri1997indexselection}) and database cracking (Idreos et~al.~\cite{idreos2007databasecracking}; Halim et~al.~\cite{halim2012stochasticcracking}) introduced techniques to adapt index layouts continuously in response to incoming queries. The Data Calculator~\cite{idreos2018datacalculator} and related synthesized key-value engines~\cite{idreos2019learningkv} advanced this concept further by modeling data structure design as a continuum of primitive structural operations synthesizable to fit arbitrary access patterns.

With the emergence of learned data structures (Kraska et~al.~\cite{kraska2018learned}), this philosophy was extended to multidimensional data. Learned multidimensional systems such as Flood (Nathan et~al.~\cite{nathan2020flood}) and Tsunami (Ding et~al.~\cite{ding2020tsunami}) train regression models to partition feature spaces dynamically according to query filter distributions, as surveyed by Pandey et~al.~\cite{pandey2020learnedspatial}. Recent workload-aware frameworks---including WAZI~\cite{pai2023wazi}, WISK~\cite{sheng2023wisk}, LiTune~\cite{wang2025litune}, ShapeShifter~\cite{wang2025shapeshifter}, and the autotuning framework of Herveau et~al.~\cite{herveau2023autotuning}---continue to refine automated partitioning and cost-model optimization for spatial and multi-attribute datasets.

However, these workload-adaptive frameworks primarily target low-dimensional range filters, point lookups, or spatial joins over tabular data, evaluating against baselines tuned specifically for relational workloads. In contrast, geometric point-cloud processing workloads feature dense, non-selective neighborhood queries (where every point queries surrounding neighbors at varying radii) over irregular three-dimensional manifolds. Fitting machine learning models across hundreds of millions of 3D coordinates incurs substantial training latency and complex spatial pruning. In Section~\ref{sec:rule}, we demonstrate that workload awareness in point clouds can be achieved with radical simplicity: a single closed-form geometric metric---the scene's global flatness ratio---reliably selects the near-optimal spatial primitive without requiring training iterations, surrogate cost models, or iterative search.

\subsection{Hardware-conscious, GPU, and out-of-core spatial search}

On massively parallel GPU architectures, spatial neighbor search is heavily constrained by thread divergence, control flow overhead, and uncoalesced memory transactions. Wald~\cite{wald2022cudakdtree} introduced cudaKDTree, which eliminates traversal stack overhead entirely through a left-balanced implicit array layout and stack-free traversal logic. Jakob and Guthe~\cite{jakob2021lbvhknn} optimized linear BVH construction and traversal specifically for GPU $k$NN queries by reorganizing point buffers for spatial locality. Taking advantage of specialized ray-tracing hardware, RTNN (Zhu~\cite{zhu2022rtnn}) reformulated nearest-neighbor search as ray-sphere intersection queries processed directly by hardware RT Cores on NVIDIA GPUs. Appendix~\ref{sec:composition} times our engine on the GPU against Wald's cudaKDTree.

When point-cloud datasets exceed physical RAM capacity, systems rely on out-of-core storage formats and streaming level-of-detail (LOD) hierarchies. Systems such as Potree (Sch\"utz et~al.~\cite{schuetz2020potree}) and the Cloud Optimized Point Cloud specification (COPC)~\cite{copc2021} store massive point collections as chunked, octree-structured files that stream pages from secondary storage or network repositories on demand. In high-resolution polygonal mesh processing, Melero et~al.~\cite{melero2019ebpoctree} employed Extended Bounding Primitive octrees (EBP-Octrees) to load whole subtrees asynchronously during collision detection.

Beyond secondary storage I/O, modern memory subsystem characteristics---notably CPU Translation Lookaside Buffers (TLB) and page table reach---exert a decisive yet frequently overlooked influence on spatial query performance. When indexing 100M to 1B points, the spatial node pool occupies hundreds of megabytes of virtual memory, triggering frequent TLB misses during pointer-chasing traversals under standard 4\,KiB virtual memory pages. By systematically evaluating page-blocked node layouts and two-megabyte Transparent Huge Pages (THP), our study bridges hardware memory architecture and spatial indexing, and Section~\ref{sec:ram} measures how much of a query's time is translation reach.

\subsection{Hybrid and nested indexes}

Because a LiDAR scene is wide and flat at the large scale and three-dimensional at the small one, several hybrid indexes stack a wide two-dimensional structure over a finer one: R$^*$-trees under an octree prefix~\cite{wang2021hybridpointcloud}, octrees nested in coarse boxes for change detection~\cite{park2021nestedoctree}, and quadtrees or grids over octrees and kd-trees for LiDAR~\cite{ogayar2023nestedlidar}, which loaded fewer nodes from an out-of-core store than either primitive alone. These comparisons were made against single structures at default parameters or under depth caps, and counted nodes rather than time. Our engine includes nesting as one option of its design space (Section~\ref{sec:designspace}); Appendix~\ref{sec:composition} reports how it fared against tuned single structures.
